\bookchapter{Capítulo 35: Crianças e celular — guia educacional na era digital}

\chapterintro{Este capítulo não trata o celular como “monstro”, nem prega que “quanto mais cedo o contato, mais inteligente”. Damos dispositivos, regras, ensinamos competências e protegemos a saúde por faixa etária, passo a passo.}

\booksection{Com que idade dar um celular à criança?}

Não há resposta padrão, mas há linhas de base por idade. Dar cedo demais — o córtex pré-frontal do cérebro infantil ainda não está maduro, e o autocontrole não acompanha o “design persuasivo” dos algoritmos. Dar tarde demais — pode causar “desconexão social” entre os colegas. Abaixo, sugestões por faixa etária:

\begin{itemize}
  \item \textbf{3-6 anos}: Evitar ao máximo o contato com celular. Se necessário, sessões únicas de no máximo 15 minutos, com conteúdo majoritariamente de animações educativas selecionadas pelos pais, e sempre\textbf{com supervisão}. Nesta fase, a tela tem o maior impacto sobre a visão e o desenvolvimento da atenção.
  \item \textbf{7-9 anos}: Pode "pegar emprestado" o dispositivo dos pais para aulas online, ouvir histórias e pesquisar. Comece a estabelecer a regra de "solicitar antes de usar, devolver após o uso", sem dispositivo próprio.
  \item \textbf{10-12 anos}: Pode-se considerar dar um "celular básico" ou relógio infantil para comunicação de ida e volta da escola. Se for realmente necessário um smartphone, ative controles parentais rigorosos.
  \item \textbf{13 anos ou mais}: No ensino fundamental II, a necessidade social aumenta. Pode-se dar um smartphone, mas com um "acordo de uso" e revisão periódica.
\end{itemize}

\begin{tcbinfo}{Princípio do "bom o suficiente"}
O primeiro celular da criança não precisa ser um flagship. O “desempenho excedente” do processador geralmente é usado para rodar jogos. Escolha um celular com bateria grande, desempenho médio e tela com proteção ocular na faixa de 1.000-2.500 RMB, ou um aparelho recondicionado oficial — isso é mais favorável ao controle. Orçamento de referência: 1.500-2.500 RMB atende à grande maioria das necessidades estudantis.

\end{tcbinfo}

\booksubsection{Relógio inteligente como solução de transição}

Se a questão é apenas “conseguir se comunicar depois da escola”, o relógio infantil é mais adequado que o celular. Ele só faz chamadas, envia mensagens de voz e localiza — sem navegador nem vídeos curtos, limitando naturalmente as “portas de entrada para o vício”:

\begin{itemize}
  \item \textbf{Xiaotiancai (Relógio Telefônico Infantil)}: Ecossistema mais completo, o "toque para adicionar amigos" entre colegas é seu principal diferencial de venda, mas o ecossistema de apps exclusivo é relativamente fechado.
  \item \textbf{Relógio Infantil Huawei}: Boa integração com celulares Huawei, localização precisa, suporte a videochamadas, bom custo-benefício.
  \item \textbf{Relógio Infantil Xiaomi}: Preço acessível, funções básicas completas, ideal para famílias com orçamento limitado.
\end{itemize}

Ao comprar um relógio, atenção: escolha modelos com suporte a\textbf{4G todas as operadoras}(2G/3G estão sendo gradualmente descontinuados) e verifique os planos infantis das operadoras.

\booksection{Guia completo de sistemas de controle parental}

Antes de entregar o dispositivo à criança, a primeira coisa não é ligar, mas configurar o controle parental. Ambas as plataformas principais têm soluções nativas:

\booksubsection{iOS: Tempo de Uso}

\begin{itemize}
  \item \textbf{Limites de Apps}: Ajustes > Tempo de Uso > Limites de Apps — permite definir tempo diário máximo por categoria de app (ex.: 30 minutos para jogos).
  \item \textbf{Tempo de Inatividade}: Defina um "horário de apagar as luzes" diário (ex.: 22h-7h), durante o qual todos os apps ficam cinza e bloqueados, apenas telefone funciona.
  \item \textbf{Restrições de Conteúdo}: Filtre filmes, sites e instalação de apps por faixa etária, e proíba compras dentro de apps.
  \item \textbf{Compartilhamento Familiar}: Use o iPhone dos pais para gerenciar remotamente o dispositivo da criança, aprovar solicitações de download e ver relatórios de uso.
\end{itemize}

\booksubsection{Android: Google Family Link / Controle parental da marca}

\begin{itemize}
  \item \textbf{Google Family Link}: Solução oficial do Google para gerenciar remotamente o tempo de uso, bloquear o dispositivo e ver relatórios de atividade. Em celulares Android na China, a disponibilidade varia conforme o modelo.
  \item \textbf{Xiaomi Espaço Infantil}: Ajustes > Espaço Infantil — permite definir tempo de uso, lista de permissões de apps; para sair, é necessária a senha dos pais.
  \item \textbf{Huawei Modo Estudante}: Ajustes > Uso saudável do celular > Modo Estudante — restringe horários e apps, com suporte a gerenciamento remoto.
  \item \textbf{OPPO/vivo Modo Infantil}: Ative em "Ajustes" ou no app exclusivo "Espaço Infantil". O princípio é semelhante.
\end{itemize}

\begin{center}
\begin{tabularx}{\linewidth}{|X|X|X|}
\hline
\textbf{Função} & \textbf{iOS Tempo de Uso} & \textbf{Android Family Link / Solução da marca} \tabularnewline
\hline
Limites de Apps & Sim, por categoria & Sim, por app \tabularnewline
Tempo de Inatividade & Sim & Sim (varia por marca) \tabularnewline
Filtro de conteúdo & Nível de sistema, mais completo & Depende da marca, qualidade variável \tabularnewline
Gerenciamento remoto & Compartilhamento Familiar & Family Link / Nuvem da marca \tabularnewline
Aprovação de instalação de apps & Requer aprovação dos pais & Family Link compatível \tabularnewline
\hline
\end{tabularx}
\end{center}

\begin{tcbwarn}{Controle parental não é "definitivo"}
As crianças sempre encontram maneiras de burlar (há inúmeros tutoriais na internet de “como quebrar o Tempo de Uso”). Medidas técnicas são apenas a linha de defesa básica — a verdadeira barreira é a comunicação e a confiança entre pais e filhos. Recomenda-se revisar mensalmente os relatórios de uso junto com a criança, em vez de monitorar unilateralmente.

\end{tcbwarn}

\booksection{Modo Adolescente}

As principais plataformas de conteúdo têm um “Modo Adolescente” integrado. Quando ativado, ele limita o tempo de uso, filtra conteúdo inadequado e desativa recargas e recompensas. Esta é uma função obrigatória por regulamentação nacional, mas\textbf{precisa ser ativada e protegida por senha pelos pais}:

\begin{itemize}
  \item \textbf{Modo Adolescente TikTok/Kuaishou}: Limitado a 40 minutos por dia (bloqueado das 22h às 6h), exibe apenas conteúdo selecionado de ciência e educação, sem transmissão ao vivo, sem recompensas, sem recargas.
  \item \textbf{Modo Adolescente WeChat}: Restringe o WeChat Pay (requer autorização de valor pelos pais), limita contas oficiais e conteúdo dos Canais, desativa funções como "Agitar" e "Pessoas Próximas" por padrão.
  \item \textbf{Sistema antiviício em jogos}: Autenticação com nome real + limite de tempo em nível nacional. Menores de idade só podem jogar às sextas, sábados, domingos e feriados legais, das 20h às 21h, com limite de 1 hora por dia.
\end{itemize}

\textbf{Não use o documento de identidade dos pais para “emprestar” autenticação}\par

Muitas crianças usam o documento de identidade dos pais para registrar contas de jogos e burlar o sistema antiviício. Isso viola os termos de uso e priva a criança de proteção. Os pais devem guardar suas senhas de pagamento e informações de identidade, e não “emprestar” a conta por conveniência.

\booksection{Educação em competências digitais}

A tecnologia pode bloquear 90\% dos riscos — os outros 10\% dependem da capacidade de julgamento da criança. Mais importante do que “proibir o acesso à internet” é ensinar “como navegar na internet”:

\begin{itemize}
  \item \textbf{Proteção de privacidade}: Não revele nome real, escola ou endereço residencial a estranhos; não envie fotos pessoais, documento de identidade ou fotos de família; mantenha a localização desativada por padrão, ativando apenas quando necessário.
  \item \textbf{Identificação e resposta ao cyberbullying}: Ser insultado, isolado ou ter fotos manipuladas maliciosamente é bullying. Ensine que, ao enfrentar isso, não revide, não apague as evidências e conte imediatamente aos pais. No WeChat, mantenha pressionada a mensagem para denunciar — a plataforma processará a conta do agressor.
  \item \textbf{Pensamento crítico sobre conteúdo}: Ensine a distinguir "fato" de "opinião", identificar clickbait e manchetes sensacionalistas; entenda que o algoritmo só mostra o que você quer ver, não a totalidade do mundo; desconfie de qualquer anúncio "bom demais para ser verdade".
  \item \textbf{Segurança online}: Não clique em links suspeitos, não escaneie QR codes de origem desconhecida, não baixe apps fora das lojas oficiais; não compartilhe senhas com ninguém (nem "melhores amigos"); considere como golpista qualquer um que peça códigos de verificação.
\end{itemize}

\begin{tcbpractice}{Mãos à obra: Assine um "Acordo de Uso do Celular" com seu filho}
Colocar as regras por escrito, com assinatura de ambos e afixado em local visível, é muito mais eficaz do que acordos verbais. Sugere-se incluir:

\begin{itemize}
  \item Horário de uso: Ex.: "Disponível das 19h30 às 20h30 após a lição de casa".
  \item Zonas proibidas: Celular guardado na mesa de jantar, no quarto antes de dormir e na sala de aula.
  \item Limites de conteúdo: Não instalar apps sem informar os pais, não conversar em privado com estranhos.
  \item Consequências por descumprimento: Ex.: "Uma hora extra, metade do tempo de uso no dia seguinte".
  \item Compromisso dos pais: Não "navegar no celular" na frente da criança por mais tempo do que o combinado, dando o exemplo.
\end{itemize}

\end{tcbpractice}

\booksection{Hábitos saudáveis de uso}

Os olhos, a coluna cervical e o sono das crianças estão em fase de desenvolvimento — o dano do celular ao corpo é maior nelas do que em adultos:

\begin{itemize}
  \item \textbf{Tempo de tela recomendado}: 3-6 anos, no máximo 1 hora por dia; 6-12 anos, sem limite para estudo, entretenimento no máximo 1 hora; acima de 12 anos, negociar com a criança, mas entretenimento no máximo 2 horas.
  \item \textbf{Hábitos de proteção ocular}: Siga a "regra 20-20-20" — a cada 20 minutos de tela, olhe para algo a 20 pés (cerca de 6 metros) de distância por 20 segundos; não use celular 1 hora antes de dormir, a luz azul inibe a melatonina e prejudica o sono.
  \item \textbf{Gestão da postura}: Tela do celular na altura dos olhos, sem cabeça baixa e ombros curvados; não use por mais de 30 minutos seguidos sem se levantar e se movimentar.
  \item \textbf{Regras conjuntas entre pais e filhos}: Estabeleça "zonas livres de celular" na mesa de jantar, no quarto e no escritório; ao assistir filmes ou fazer atividades ao ar livre em família, guarde todos os celulares — faça do "largar o celular" uma cultura familiar, não uma punição para a criança.
\end{itemize}

\begin{tcbinfo}{Os pais são o melhor "firewall"}
Pesquisas mostram que o tempo de uso do celular pelos pais tem forte correlação positiva com o dos filhos. Exigir que a criança não use o celular enquanto você mesmo navega em vídeos curtos até tarde da noite torna as regras ineficazes. Para mudar a criança, comece largando seu próprio celular.

\end{tcbinfo}

\booksection{Herança digital e crianças}

Esta seção é facilmente negligenciada, mas é muito importante. Os ativos digitais das crianças (contas, fotos, saldo de recargas, equipamentos de jogos) também enfrentam questões de herança e privacidade:

\begin{itemize}
  \item \textbf{Herança de ativos digitais infantis}: Em caso de falecimento da criança, o processo de herança de suas contas de redes sociais, jogos e fotos na nuvem é igual ao de adultos — requer certidão de óbito e comprovante de parentesco. Recomenda-se que os pais configurem desde cedo um "contato herdeiro" ou "contato de emergência" para as contas importantes da criança (como contato herdeiro do Apple, herança de conta do WeChat).
  \item \textbf{Privacidade digital de menores}: A Lei de Proteção de Informações Pessoais determina que o tratamento de informações pessoais de menores de 14 anos requer consentimento do responsável legal e deve ter regras de privacidade específicas. Os pais devem revisar regularmente as políticas de privacidade dos apps usados pela criança e desativar as opções de "recomendações personalizadas" e "compartilhamento de dados".
  \item \textbf{Sugestão para registro de contas}: As contas criadas para a criança devem usar e-mail/número de telefone controlado pelos pais, facilitando a gestão unificada durante a menoridade; quando a criança atingir a maioridade, a conta pode ser transferida.
\end{itemize}

\begin{tcbwarn}{Alerta sobre o outro lado da "herança digital infantil"}
As crianças vêm sendo “registradas” desde o nascimento — fotos de crescimento postadas pelos pais, vídeos enviados pela escola, dados de saúde do relógio inteligente. Uma vez carregados, esses dados são difíceis de excluir completamente. Antes de publicar, pense: daqui a dez anos, a criança ficaria constrangida com esta foto? Respeitar a personalidade digital da criança começa com a moderação ao “expor os filhos”.

\end{tcbwarn}

\begin{tcbsummary}{Resumo do capítulo}
\begin{itemize}
  \item Não há idade padrão para dar o celular, mas evite ao máximo o contato entre 3-6 anos; antes dos 13, use relógio infantil ou celular básico como transição
  \item "Bom o suficiente": o primeiro celular deve ser um modelo de entrada, não flagship — desempenho excedente equivale a tentação excedente
  \item iOS Tempo de Uso e Android Family Link/controle parental da marca são duas linhas de defesa técnicas — use-as em conjunto para maior segurança
  \item Modo Adolescente e sistema antiviício em jogos são exigências legais, mas precisam ser ativados e protegidos por senha pelos pais
  \item Competências digitais são mais importantes que bloqueio técnico: ensine a identificar bullying, proteger privacidade e pensar criticamente sobre o conteúdo
  \item Hábitos saudáveis começam pelos pais — quando você larga o celular, a criança também larga
  \item Herança digital infantil e privacidade são temas emergentes — configure o "contato herdeiro" desde cedo e modere a exposição dos filhos
\end{itemize}

\end{tcbsummary}
